\begin{abstract}
   The 2020s is the decade of survey instruments in astronomy.
   Radio astronomy is no exception, with Caltech's proposed DSA-2000 being the most powerful radio interferometer in the world, costing much less than competing instruments.
   Key to this achievement are two core breakthroughs, a completely ambient-temperature receiver and a ``radio camera'' backend that images the sky in real time.
   DSA-2000 will have record-breaking survey speed and sensitivity, enabled by these two key breakthroughs, giving astronomers all over the world open access to exquisite all-sky maps to enable the discovery of billions of new radio sources, precise timing of pulsars, and localization of fast radio bursts.
   The array will produce enough data to keep astronomers busy for a century.

   In this thesis, we discuss the development of one of the key breakthroughs, the ambient-temperature receiver.
   Specifically, we focus on the design, testing, and implementation of the wideband, ambient-temperature low noise amplifier.
   We cover the design from analytic first principles through precision measurement of its performance.
   We follow this with a discussion of the design and implementation of the analog signal path, including a high performance, analog over fiber backhaul link.
   Finally, we discuss the Galactic Radio Explorer (GReX) instrument, designed as a global experiment probing the brightest radio transients in the local universe.
\end{abstract}